% =====================================================================
% 卒業研究1 中間発表 概要（2ページ）
% 提出締切: 2026-07-17 12:00 LETUS
% 著者: 指田一茶（東京理科大学 創域理工学部 経営システム工学科 秦野研究室）
% バージョン: v1（友人フォーマット・執筆要項準拠版 2026-07-13）
% =====================================================================
\documentclass[10pt,a4paper,twocolumn]{article}

% --- フォント・日本語 ---
\usepackage[T1]{fontenc}
\usepackage{lmodern}
% 日本語環境では以下を有効化（upLaTeX / LuaLaTeX 使用時）
% \usepackage{luatexja}

% --- レイアウト ---
\usepackage[top=20mm,bottom=20mm,left=15mm,right=15mm,columnsep=8mm]{geometry}
\usepackage{setspace}
\setstretch{1.1}

% --- 数式・記号 ---
\usepackage{amsmath,amssymb}

% --- 参考文献 ---
\usepackage{cite}

% --- タイトル余白の削減 ---
\usepackage{titling}
\setlength{\droptitle}{-2em}

% =====================================================================
\title{\large\bfseries Murphy 分解を用いた LLM 確信度プロンプト設計の機序解明\\
{\normalsize --- 日本語タスクにおける Calibration (REL) と Resolution (RES) の分離分析 ---}}

\author{\normalsize 秦野研究室 \quad 7423XXX \quad 指田 一茶}

\date{}
% =====================================================================

\begin{document}
\maketitle
\thispagestyle{empty}

% =====================================================================
\section{研究背景}
% =====================================================================

大規模言語モデル（LLM）は教育・業務・日常生活へ急速に普及しているが，
「ハルシネーション」，すなわち誤情報を自信満々に生成する問題は未解決である\cite{ji2023survey}。
問題の本質は，モデルが示す確信度と実際の正答率の乖離にあり，
確信度が正答率と整合していれば（well-calibrated），利用者はそれをリスク判断の手がかりにできる。

キャリブレーション改善手法は大別して Post-hoc 補正（Temperature Scaling 等），
ファインチューニング，プロンプト設計の3系統がある。
GPT・Claude・Gemini のような商用 API は内部 logits やパラメータへのアクセスを提供しないため，
プロンプト設計のみが唯一実践可能な手段となる。
先行研究\cite{tian2023just,xiong2024can}は，
英語タスクにおいて「言語化された確信度（verbalized confidence）」を
プロンプトで引き出す方式が内部 softmax 確率より優れた較正精度を達成することを示した。
しかし，これらの研究はプロンプト方式の効果を ECE や Brier Score といった単一の総合指標でのみ評価しており，
\textbf{なぜ}特定のプロンプト設計が有効かという機序を明らかにしていない。
Murphy\cite{murphy1973vector}は Brier Score を
$\text{BS} = \text{REL} - \text{RES} + \text{UNC}$
の3成分に分解できることを示し，Bröcker\cite{brocker2009reliability}はこの分解が
任意の Proper Scoring Rule に対して成立することを一般化した。
著者らが調査した範囲では，この Murphy 分解を LLM の確信度評価に適用した研究は現時点で存在しない。

% =====================================================================
\section{研究目的}
% =====================================================================

本研究は Murphy 分解を LLM 評価に初めて適用し，
プロンプト設計の効果が\textbf{較正精度（REL）の改善}に由来するのか，
\textbf{識別能力（RES）の改善}に由来するのかを定量的に解明することを目的とする。

\textbf{メイン RQ}：確信度プロンプト設計の効果は，REL の改善か，RES の改善か，
それとも両者の複合か？

\textbf{サブ RQ}：モデル間・タスクドメイン間・プロンプト方式間・日英言語間で
REL・RES の分離パターンにどのような差があるか？

演習2では ECE を単一指標としたプロンプト間比較にとどまったが，
本研究はその効果の内訳（REL/RES）を分離することで機序に踏み込む点で
演習2の焼き直しではなく，卒研2で本格実施する本実験・機序解明の考察へ直結する。

% =====================================================================
\section{研究の方針}
% =====================================================================

\textbf{モデル}：GPT-4o，Claude Sonnet 4.6，Gemini 2.x，Swallow（4〜6 モデル）。

\textbf{データセット}：4 ドメイン × 250 問 = 計 1,000 問
（MGSM・JCommonsenseQA・JMMLU・FLORES-200）。

\textbf{プロンプト方式}（3 方式を比較）：
Verb.1S（回答と確信度を同時出力），
Verb.2S（Turn1 で回答，Turn2 で確信度を別途質問），
Ling.1S（言語表現5段階を数値マッピング）。
いずれもプロンプト設計のみで実現し，logits・再学習は不要である。

\textbf{評価指標}：

\begin{align}
  \text{ECE} &= \sum_{m=1}^{M} \frac{|B_m|}{n} \left|\text{acc}(B_m) - \text{conf}(B_m)\right| \\
  \text{BS} &= \frac{1}{N} \sum_{i=1}^{N} (p_i - y_i)^2
\end{align}

Murphy 分解：$\text{BS} = \text{REL} - \text{RES} + \text{UNC}$，
$\text{REL} = \sum_k \frac{n_k}{n}(\bar{f}_k - \bar{o}_k)^2$（低いほど良），
$\text{RES} = \sum_k \frac{n_k}{n}(\bar{o}_k - \bar{o})^2$（高いほど良）。
補助指標として AUROC（識別能力）と
Ferro \& Fricker 分散推定\cite{ferro2012decomposition}による信頼区間を用いる。

\textbf{正答判定}：数学（MGSM）は正規化完全一致，
常識・知識（JCommonsenseQA・JMMLU）は選択肢ラベル一致，
翻訳（FLORES-200）は COMET スコア閾値（50文の人間判定で F1 最適化）とする。

\textbf{検証仮説}：H1（Verb.2S は RES を改善），
H2（Verb.1S は REL を改善），
H3（Ling.1S は両指標とも改善が弱い），
H4（日本語は英語より REL の改善幅が小さい）。

% =====================================================================
\section{現在の進捗}
% =====================================================================

評価指標・データセット・プロンプト設計は2026年6月のゼミで承認済みである。
Murphy 分解の計算コードを実装し，完全較正な合成データで
$\text{BS} = \text{REL} - \text{RES} + \text{UNC}$ の等式を数値的に検証した
（誤差は分位点分割に起因する既知の範囲内）。
3方式（Verb.1S/Verb.2S/Ling.1S）に対応したプロンプトテンプレートと
実験実行スクリプトの整備を完了し，横断調査により
LLM 論文における Murphy 分解の適用例が現時点で存在しないことを確認した。
一方，API を用いた本実験・パイロット実験は未実施であり，
実測データに基づく REL/RES の分離結果はまだ得られていない。

% =====================================================================
\section{今後の予定}
% =====================================================================

2026年7月中旬にパイロット実験（Claude 1モデル・少数問題）でスクリプトの動作確認を行い，
問題がなければ7月末までに本実験の一部を開始する。
2026年8〜9月にかけて4モデル×3方式×1,000問の本実験を実施し，REL/RES の分離分析を行う。
2026年10〜11月に機序解明の考察を中心とした各章の草稿を作成し，
2026年12月に卒研2としての最終稿をまとめ，2027年1月に卒業論文として提出する予定である。
本研究は演習2で確立した評価基盤を土台としつつ，機序解明という新たな軸を加えることで
卒研2の本実験・考察に直接接続する構成となっている。

% =====================================================================
\section*{参考文献}
% =====================================================================

\begin{thebibliography}{9}

\bibitem{ji2023survey}
Ji, Z., et al. (2023).
Survey of Hallucination in Natural Language Generation.
\textit{ACM Computing Surveys}, 55(12), 1--38.

\bibitem{tian2023just}
Tian, K., et al. (2023).
Just Ask for Calibration: Strategies for Eliciting Calibrated Confidence Scores
from Language Models Fine-Tuned with Human Feedback.
\textit{Proc. EMNLP 2023}.

\bibitem{xiong2024can}
Xiong, M., et al. (2024).
Can LLMs Express Their Uncertainty? An Empirical Evaluation of Confidence
Elicitation in LLMs.
\textit{Proc. ICLR 2024}.

\bibitem{murphy1973vector}
Murphy, A. H. (1973).
A New Vector Partition of the Probability Score.
\textit{J. Applied Meteorology}, 12(4), 595--600.

\bibitem{brocker2009reliability}
Bröcker, J. (2009).
Reliability, Sufficiency, and the Decomposition of Proper Scores.
\textit{QJRMS}, 135(643), 1512--1519.

\bibitem{ferro2012decomposition}
Ferro, C. A. T., \& Fricker, T. E. (2012).
A Bias-Corrected Decomposition of the Brier Score.
\textit{QJRMS}, 138(668), 1954--1960.

\bibitem{guo2017calibration}
Guo, C., et al. (2017).
On Calibration of Modern Neural Networks.
\textit{Proc. ICML 2017}.

\end{thebibliography}

\end{document}
